%\begin{center}
%\large \bf \runtitulo
%\end{center}
%\vspace{1cm}
\chapter*{\runtitulo}

\noindent La resistencia a los antimicrobianos en \textit{Acinetobacter baumannii} constituye un desafío clínico prioritario debido a la frecuencia de infecciones asociadas al cuidado de la salud y a la limitada disponibilidad de opciones terapéuticas eficaces. La identificación de nuevas dianas antibacterianas requiere integrar información que refleje tanto su relevancia funcional como su conservación y su potencial de intervención farmacológica. Esta tesis desarrolla un flujo computacional multicriterio para priorizar potenciales dianas terapéuticas en cepas clínicas de \textit{A.~baumannii}. El enfoque integra el análisis de pangenoma para identificar genes conservados, evidencia experimental publicada de esencialidad genética, anotación funcional, participación en rutas metabólicas y evaluación de homología con proteínas humanas. Los candidatos mejor posicionados se caracterizan además mediante el análisis de estructuras proteicas y cavidades de unión, con el fin de estimar su ligabilidad. La integración de estas capas de evidencia permite generar un ranking reproducible, interpretable y trazable, que explicita las fortalezas y limitaciones de cada candidato. El resultado del trabajo constituye una estrategia para reducir el espacio de búsqueda de dianas y orientar la posterior validación experimental y el desarrollo racional de nuevos antibacterianos contra \textit{A.~baumannii}.

\bigskip

\noindent\textbf{Palabras clave:} \textit{Acinetobacter baumannii}, resistencia a los antimicrobianos, dianas terapéuticas, esencialidad genética, pangenoma, ligabilidad.
